\section{The finite classification}\label{sec:finite-statement}
A \emph{Grothendieck ring} here means a fusion ring admitting a fusion-category
categorification over $\mathbb C$. The question concerns existence for a
based ring, not a classification of all tensor categories realizing it.
For structure constants $N_{ij}^k$, put
$\mu(R)=\max_{i,j,k}N_{ij}^k$. We use the pinned rank-four census
\cite{Census} and write $F_{m,j}$ for its $j$th ring of exact multiplicity
$m$. For compatibility with historical lists, the identifiers through
$m=32$ follow the mapping supplied with the certificates; above $32$ they
follow the pinned parameter-file order. The supplement records the source
index, tensor digest, and any historical basis permutation. Identifiers
name based rings but are never premises of an obstruction test.

\begin{theorem}[Bounded classification]\label{thm:finite}
A rank-four fusion ring $R$ with $\mu(R)\leq\MultiplicityBound$ admits a fusion
categorification over $\mathbb C$ if and only if it appears in
Table~\ref{tab:models}. Every ring in that table admits a unitary
categorification. Consequently general, pivotal, spherical, pseudo-unitary
and unitary categorifiability give the same fifteen based rings in this range.
\end{theorem}

The forward implication is the exclusion theorem proved below; the reverse
implication uses the established models in the table. The logical reduction
is particularly useful: a unitary realization establishes all five existence
statements, whereas a general obstruction excludes all five. There is no
need for five independent classification algorithms.

\begin{table}[ht]
\centering
\begin{tabular}{lll}
\toprule
Ring & A unitary realization & Source for existence\\
\midrule
$F_{1,1}$ & $\Vec_{C_2\times C_2}$ & \cite{LPR}\\
$F_{1,2}$ & $\mathrm{Fib}\boxtimes\Vec_{C_2}$ & \cite{LPR}\\
$F_{1,3}$ & $\Rep(D_{10})$ & \cite{LPR}\\
$F_{1,4}$ & $\mathrm{PSU}(2)_6$ & \cite{LPR}\\
$F_{1,5}$ & $\mathrm{Fib}\boxtimes\mathrm{Fib}$ & \cite{LPR}\\
$F_{1,6}$ & $\mathrm{PSU}(2)_7$ & \cite{LPR}\\
$F_{1,7}$ & $\Vec_{C_4}$ & \cite{LPR}\\
$F_{1,8}$ & $\mathrm{TY}(C_3)$ & \cite{LPR}\\
$F_{1,10}$ & the pseudo-$\mathrm{PSU}(2)_6$ ring & \cite{LPR}\\
$F_{2,2}$ & the even part of $2D2$ & \cite[Theorem 3.1]{EIP}\\
$F_{2,15}$ & the Haagerup category $\mathcal H_1$ & \cite[Table 2]{GrossmanSnyder}\\
$F_{2,17}$ & $\Rep(A_4)$ & standard representation category\\
$F_{3,19}$ & $\mathrm{Ad}(E_8)$ & \cite{EvansJones}\\
$F_{3,24}$ & $C_3$ near-group, parameter $3$ & \cite[Theorem 9.1]{Izumi}\\
$F_{6,81}$ & $C_3$ near-group, parameter $6$ & \cite[Theorem 10.19]{Izumi}\\
\bottomrule
\end{tabular}
\caption{The fifteen positive rings. Here $D_{10}$ has order ten;
$\mathrm{Fib}$ is unitary Fibonacci, and $\mathrm{PSU}(2)_k$ denotes the even
part of the standard unitary $\mathrm{SU}(2)_k$ category. The word
``pseudo'' in the name of $F_{1,10}$ does not weaken its unitary realization.}
\label{tab:models}
\end{table}

The census contains $\CensusSize$ based rings through multiplicity
$\MultiplicityBound$. Table~\ref{tab:census-bands} groups the exact counts
in four disjoint ranges; the supplement gives the count at every exact
multiplicity and a split by duality type.
% Generated from audited census and verified certificate summary.
\begin{table}[ht]
\centering
\begin{tabular}{lrr}
\toprule
Exact multiplicity range & Fusion rings & Categorifiable rings\\
\midrule
$1$--$32$ & $17\,736$ & $15$\\
$33$--$64$ & $104\,672$ & $0$\\
$65$--$96$ & $268\,810$ & $0$\\
$97$--$128$ & $498\,312$ & $0$\\
\midrule
Total & $\CensusSize$ & $15$\\
\bottomrule
\end{tabular}
\caption{The complete census and its categorifiable subset.}
\label{tab:census-bands}
\end{table}

The input is the rank-four data of the fusion-ring census
\cite{Census} at commit
\texttt{7f1b36e74df53a8a1cf8597db45d261e408d61b2}.
That repository supplies exhaustive coverage and exact counts through
$m=128$; Appendix~\ref{sec:census} records the tensor interface needed by
the classification checker. Every identifier is accompanied by a full
tensor digest, so no mathematical conclusion depends on catalogue order.
Every input ring is covered by a realization or nonexistence certificate.

The mathematical tools are presented before their finite application. Known
results are cited with their hypotheses; the new inequalities and reductions
are proved. The electronic supplement records applications as coefficient,
polynomial, prime, sign and bounded-enumeration certificates, rather than as
hundreds of separate case proofs. The companion article proves the
unrestricted rank-four classification and automatic sphericality. The
finite proof of Theorem~\ref{thm:finite} is independent of that article.
